\chapter{Handling Mathematical Expression Data}\label{mexpr}

This chapter introduces Egison's features for handling symbolic computation.

\section{Internal Representation of Mathematical Expression Data}\label{mexpr-expr}

Mathematical expressions are directly implemented as a built-in data type within the Egison interpreter using Haskell's algebraic data types.
By using the built-in function \lstinline{fromMathExpr}, this mathematical expression data can be converted into an Egison algebraic data type.
\begin{lstlisting}[language=egison,numbers=none]
fromMathExpr  (x - sqrt 2)
-- Div (Plus [Term 1 [(Symbol x [], 1)], Term -1 [(Apply sqrt [Div (Plus [Term 2 []])
--     (Plus [Term 1 []])], 1)]]) (Plus [Term 1 []])
\end{lstlisting}
This section explains the internal algebraic data type that represents mathematical expression data.

As can be seen from the algebraic data type in the output above, the mathematical expression data type is defined as a fraction that takes polynomials in sum-of-products normal form for both the numerator and denominator.
\lstinline{Div} is the data constructor for creating this fraction, and \lstinline{Plus} is the data constructor for creating polynomials.
Therefore, for example, the result of adding two fractions is combined into a single fraction.
\begin{lstlisting}[language=egison,numbers=none]
a / b + x / y
-- (a * y + x * b) / (b * y)
\end{lstlisting}
As can be seen from the fact that the argument of \lstinline{Plus} is a list, a polynomial is defined as a data type consisting of a list of terms.
A term is defined as a data type consisting of a coefficient and a list of factors.
Terms are created using the \lstinline{Term} data constructor.
\lstinline{Term} takes a coefficient as the first argument and a list of factors as the second.
The coefficient can be an integer, and factors can be symbols, among other things.
For example, the term \lstinline{3 * x^2 * sqrt 2} has \lstinline{3} as the coefficient and \lstinline{[x^2, sqrt 2]} as the list of factors.
Factors can include not only symbols but also terms like $\sqrt{2}$.
The list of factors is an association list of bases (e.g., $x$ for $x^2$) and exponents (e.g., $2$ for $x^2$).
There is a restriction that exponents can only be natural numbers.
\begin{lstlisting}[language=egison,numbers=none]
fromMathExpr (3 * x^2 * sqrt 2)
-- Div (Plus [Term 3 [(Symbol x [], 2), (Apply sqrt [Div (Plus [Term 2 []]) (Plus [Term 1 []])], 1)]]) (Plus [Term 1 []])
\end{lstlisting}
As shown in the example above, factors are created using data constructors such as \lstinline{Symbol} and \lstinline{Apply}.
The \lstinline{Symbol} data constructor takes the symbol itself as the first argument and a list of indices as the second.
The reason for placing the symbol itself as the first argument of the \lstinline{Symbol} data constructor is to make the result of \lstinline{fromMathExpr} easier to work with.
Since having the symbol itself within the data structure representing a symbol would create an infinite circular definition, the internal Haskell representation of symbols uses a string as the first argument of the \lstinline{Symbol} data constructor.
Indices are explained in Chapter~\ref{tensor}, which covers tensors.
The \lstinline{Apply} data constructor takes a function name symbol as the first argument and a list of argument expressions as the second.
There are also \lstinline{Quote} and \lstinline{Function} as factor data constructors.
The \lstinline{Quote} data constructor is introduced in Section~\ref{mexpr-backquote}.
The \lstinline{Function} data constructor is introduced in Chapter~\ref{funcsym}.

A \lstinline{toMathExpr} function, which is the inverse of \lstinline{fromMathExpr}, is also provided to convert mathematical expressions represented as Egison algebraic data types back into Egison's internal mathematical expression data.
Using \lstinline{toMathExpr}, after processing mathematical expressions as Egison algebraic data types, the data can be converted back to Egison's internal mathematical expression data.
\begin{lstlisting}[language=egison,numbers=none]
toMathExpr (fromMathExpr  (x - sqrt 2))
-- x - sqrt 2
\end{lstlisting}
However, built-in functions like \lstinline{fromMathExpr} and \lstinline{toMathExpr} rarely need to be used directly by the user, as the \lstinline{mathExpr} matcher defined in the Egison library can be used instead.
\lstinline{fromMathExpr} and \lstinline{toMathExpr} are used in the implementation of the \lstinline{mathExpr} matcher.

Numbers are a complex data type.
Up to natural numbers, they can be elegantly defined as algebraic data types.
However, when considering subtraction, the inverse of addition over natural numbers, the concept of negative numbers emerges.
Furthermore, considering the inverse of multiplication, which is obtained by repeated addition, leads to the concept of rational numbers.
Moreover, considering the inverse of exponentiation, which is obtained by repeated multiplication, leads to the concepts of square roots, imaginary numbers, and algebraic numbers that include them.

\section{Controlling Expression Expansion with Backquote (\texttt{`})}\label{mexpr-backquote}

As mentioned in Section~\ref{cas-symbol}, Egison automatically expands mathematical expressions into sum-of-products normal form, but this expansion can be controlled using the backquote (\lstinline{`}).
An expression following a backquote is treated as a single symbol.
\begin{lstlisting}[language=egison,numbers=none,escapechar=']
`(x + 1)^2
--  `(x + 1)^2

`(x + 1) + `(x + 1)
--  2 * `(x + 1)

`(x + 1) + (x + 1)
-- `(x + 1) + x + 1
\end{lstlisting}

Backquotes are mainly useful in computations involving polynomial division, as shown below.
\begin{lstlisting}[language=egison,numbers=none,escapechar=']
(x + 1)^2 / (x + 1)
-- (x^2 + 2 * x + 1) / (x + 1)

`(x + 1)^2 / `(x + 1)
-- '(x + 1)
\end{lstlisting}

A quoted term is treated as a factor within Egison's internal mathematical expression data.
\lstinline{Quote} is defined as a data constructor at the same level as \lstinline{Symbol} and \lstinline{Apply}.
\begin{lstlisting}[language=egison,numbers=none,escapechar=']
fromMathExpr (`(x + 1)^2)
-- Div (Plus [Term 1 [(Quote (Div (Plus [Term 1 [(Symbol x [], 1)], Term 1 []]) (Plus [Term 1 []])), 2)]]) (Plus [Term 1 []])fromMathExpr `(x + 1)^2
\end{lstlisting}

\section{Controlling Function Application with Single Quote (\texttt{'})}\label{mexpr-squote}

Egison returns a symbolic result when a symbol is applied as a function.
For example, function application of an undefined variable \lstinline{f} behaves as follows.
\begin{lstlisting}[language=egison,numbers=none]
f a -- f a
\end{lstlisting}
Sometimes we want defined functions to also return symbolic results as above.
For example, for the \lstinline{sqrt} function, we want the following behavior.
\begin{lstlisting}[language=egison,numbers=none]
sqrt 4 -- 2
sqrt 5 -- sqrt 5
\end{lstlisting}
To achieve this behavior, Egison provides the single quote (\lstinline{'}) as a built-in syntax.
A single quote is followed by a variable, and it returns the symbol with that variable's name, even if the variable is already defined.
By prefixing a function with a single quote, symbolic function application can be expressed even for defined functions.
\begin{lstlisting}[language=egison,numbers=none]
def f (x: MathExpr) : MathExpr := x + 1

f a -- a + 1
'f a -- f a
\end{lstlisting}
Beyond the implementation of \lstinline{sqrt} for computing square roots, single quotes are also used in the Egison library for implementing trigonometric functions such as \lstinline{sin} and \lstinline{cos}.

\begin{lstlisting}[language=egison,numbers=none]
fromMathExpr x
-- Div (Plus [Term 1 [(Symbol x [], 1)]]) (Plus [Term 1 []])
fromMathExpr 'x
-- Div (Plus [Term 1 [(Symbol x [], 1)]]) (Plus [Term 1 []])
\end{lstlisting}

\section{Declaring Local Symbols with \texttt{withSymbols} Expressions}\label{mexpr-withsym}

The specification that undefined variables are treated as symbols tends to lead to mistakes where users attempt to use already-defined variables as symbols.
To solve this problem, Egison provides the \lstinline{withSymbols} expression syntax for declaring local symbols.
The \lstinline{withSymbols} expression takes a list of variables as its first argument, and these variables are treated as symbols during the evaluation of the expression in the second argument.
\begin{lstlisting}[language=egison,numbers=none]
def i : Integer := 10

i -- 10
withSymbols [i] i -- i
i + withSymbols [i] i -- 10 + i
\end{lstlisting}

Local symbols declared with \lstinline{withSymbols} are treated as independent symbols from global symbols generated from undefined variables.
\begin{lstlisting}[language=egison,numbers=none]
j + withSymbols [j] j -- j + j
\end{lstlisting}

\section{Pattern Matching on Mathematical Expression Data}\label{mexpr-pattern}

This section introduces the patterns available for mathematical expression data.
The matcher \lstinline{mathExpr} for mathematical expressions is implemented in the Egison library.
The definition of \lstinline{mathExpr} can be found in \texttt{lib/math/expression.egi}.
Additionally, several syntactic sugar constructs are built in to allow patterns on mathematical expressions to be written in a form closer to mathematical notation.
This section also introduces these syntactic sugar constructs.
Pattern matching on mathematical expressions is demonstrated in Section~\ref{cas-derivative}.
Please refer to the program in Section~\ref{cas-derivative} for concrete examples of pattern usage.

\subsection{Pattern Matching on Factors}

There are four pattern constructors for factors: \lstinline{symbol}, \lstinline{apply}, \lstinline{quote}, and \lstinline{func}.
\lstinline{symbol} matches symbols, \lstinline{apply} matches symbolic function applications introduced in Section~\ref{cas-symbol2}, \lstinline{quote} matches backquoted expressions introduced in Section~\ref{mexpr-backquote}, and \lstinline{func} matches objects called function symbols, which are explained in Chapter~\ref{funcsym}.

The first argument of the \lstinline{symbol} pattern constructor matches the symbol itself, and the second argument matches the indices.
Indices are discussed again in Chapter~\ref{tensor}, which covers tensors.
\begin{lstlisting}[language=egison,numbers=none]
match x as mathExpr with
  | symbol $s _ -> s
-- x
\end{lstlisting}
While one can use the \lstinline{symbol} pattern constructor to match symbols, in practice, value patterns are more commonly used.
\begin{lstlisting}[language=egison,numbers=none]
match x as mathExpr with
  | #x -> "Matched"
  | _ -> "Not matched"
-- "Matched"
\end{lstlisting}

The first argument of the \lstinline{apply} pattern constructor matches the function name symbol, and the second argument matches the list of argument expressions.
Since the \lstinline{list mathExpr} matcher is used for matching the arguments, list patterns can be used.
\begin{lstlisting}[language=egison,numbers=none]
match ('f x 1 y)  as mathExpr with
  | apply $f ($x :: $xs) -> (f, x, xs)
-- (f, x, [1, y])
\end{lstlisting}

The \lstinline{quote} pattern constructor takes a pattern for a mathematical expression as its argument.
This pattern is matched against the contents of the quote using the \lstinline{mathExpr} matcher.
\begin{lstlisting}[language=egison,numbers=none,escapechar=']
match `(x + 1) as mathExpr with
  | quote $e -> e
-- x + 1
\end{lstlisting}


\subsection{Pattern Matching on Terms}

The \lstinline{term} pattern constructor is used for pattern matching on terms.
The first argument of the \lstinline{term} pattern constructor matches the coefficient, and the second argument matches the association list of factors.
The assocMultiset matcher introduced in Section~\ref{matcher-assoc} is used for pattern matching on the association list of factors.
\begin{lstlisting}[language=egison,numbers=none]
matchAll 3 * x^2 * y  as mathExpr with
  | term $a ($x :: $xs) -> (a, x, xs)
-- [(3, x, [(x, 1), (y, 1)]), (3, y, [(x, 2)])]
\end{lstlisting}

There is syntactic sugar for the \lstinline{term} pattern constructor that allows writing patterns closer to mathematical notation.
\begin{lstlisting}[language=egison,numbers=none]
matchAll 3 * x^2 * y  as mathExpr with
  | $a * $x * $xs -> (a, x, xs)
-- [(3, x, x * y), (3, y, x^2)]

matchAll 3 * x^2 * y  as mathExpr with
  | $a * $x^#2 * $xs -> (a, x, xs)
-- [(3, x, y)]
\end{lstlisting}

\subsection{Pattern Matching on Polynomials}

The \lstinline{poly} pattern constructor is used for pattern matching on rational expressions.
The \lstinline{poly} pattern constructor takes a pattern that matches a list of terms as its argument.
\begin{lstlisting}[language=egison,numbers=none]
matchAll x + y + 1  as mathExpr with
  | poly ($x :: $xs) -> (x, xs)
-- [(x, [y, 1]), (y, [x, 1]), (1, [x, y])]
\end{lstlisting}

There is also syntactic sugar for the \lstinline{poly} pattern constructor to allow writing patterns closer to mathematical notation.
\begin{lstlisting}[language=egison,numbers=none]
matchAll x + y + 1  as mathExpr with
  | $x + $xs -> (x, xs)
-- [(x, y + 1), (y, x + 1), (1, x + y)]
\end{lstlisting}

\subsection{Pattern Matching on Rational Expressions}

The \lstinline{div} pattern constructor is used for pattern matching on rational expressions.
The first argument of the \lstinline{div} pattern constructor matches the numerator polynomial, and the second argument matches the denominator polynomial.
\begin{lstlisting}[language=egison,numbers=none]
matchAll (a + b) / c  as mathExpr with
  | div $x $y -> (x, y)
-- [(a + b, c)]
\end{lstlisting}

There is also syntactic sugar for the \lstinline{div} pattern constructor to allow writing patterns closer to mathematical notation.
\begin{lstlisting}[language=egison,numbers=none]
matchAll (a + b) / c  as mathExpr with
  | $x / $y -> (x, y)
-- [(a + b, c)]
\end{lstlisting}
